\section{Introduction}
\label{sec:intro}

\begin{itemize}
    \item Who is the organisation
    \item What does it do?
    \item why is it relevant?
    \item problem addressed
    \item general inforamtion about the company
\end{itemize}

Läderach is a Swiss premium chocolate manufacturer and retailer headquartered in
Ennenda, Canton of Glarus, Switzerland. Founded in 1962, the company has developed
from a specialised supplier of Swiss chocolate products into an internationally
active confectionery brand with its own boutiques, franchise partnerships, and a
growing global retail presence \cite{laderachOfficialWebsite}. Its main operating
entity, Läderach (Schweiz) AG, is organised as an Aktiengesellschaft under Swiss
law and is registered at Bleiche 14, 8755 Ennenda, Switzerland
\cite{moneyhouseLaederachSchweizAG}. The company describes itself as fully
family-owned, with members of the Läderach family continuing to play central roles
in ownership, governance, and executive management \cite{laderachLeadershipNextGeneration}.

The organisation operates in the premium chocolate and confectionery industry. Its
core activities include cocoa sourcing, chocolate mass production, product
manufacturing, supply-chain management, and retail distribution through branded
chocolateries \cite{laderachOurBusinessUnits}. Läderach emphasises a vertically
integrated business model, often described as control from the cocoa bean to the
shop counter, and states that its chocolate production is carried out in
Switzerland \cite{laderachBestManaged2025}. This integration is relevant because
quality, freshness, traceability, and brand experience are central sources of
competitive advantage in the premium chocolate market.

Läderach is also relevant as an example of a privately held Swiss family business
that has pursued rapid international expansion while maintaining a strong
association with Swiss craftsmanship and premium positioning. In 2024, the company
announced the opening of its 200th chocolaterie in London and reported continued
sales growth, particularly driven by international locations
\cite{laderach200thChocolaterie2024}. By 2025, Läderach reported more than 2,500
employees from over 80 nationalities, indicating both organisational growth and an
increasingly international workforce \cite{laderachBestManaged2025}.

From an analytical perspective, Läderach provides a useful case for examining the
tensions faced by premium consumer-goods companies: the need to scale globally
while preserving perceived authenticity, product quality, supply-chain credibility,
and reputational trust. Although specific challenges will be discussed in greater
detail later, the company’s development raises broader questions about how
family-owned luxury food brands manage growth, governance, stakeholder
expectations, and public reputation in an increasingly transparent and competitive
market environment.